% English abstract (11pt Shinmyungjo style)

\noindent
This thesis presents \textbf{QASAMAP} --- the Quantum-Agentic Smart Manufacturing Anomaly detection and Predictive Maintenance framework --- structured as a three-layer real-time pipeline for industrial predictive maintenance. The framework is grounded on a 100{,}000-reading smart manufacturing dataset spanning 50 machines over a 69-day monitoring period, with five sensor channels (temperature, vibration, humidity, pressure, energy consumption), five failure modes, and a streaming infrastructure that delivers data every minute via Apache Kafka.

\textbf{Layer~1} performs spatiotemporal sensor forecasting through a Temporal Graph Convolutional Network (T-GCN) constructed over a directed cross-correlation lag graph (50 nodes, 538 edges, Machine~40 anchor). The T-GCN achieves average R\textsuperscript{2} = 0.7702 across five sensor channels, outperforming both standalone LSTM (R\textsuperscript{2} = 0.6853) and Hybrid LSTM-GCN (R\textsuperscript{2} = 0.7134). Supervised classifiers (XGBoost, LightGBM, CatBoost, HistGradientBoosting, TabNet) provide companion outputs (anomaly flag, downtime risk, RUL forecast, failure type, maintenance flag) with best F1 of 0.999 on anomaly detection and 0.880 on five-class failure type classification.

\textbf{Layer~2} implements a three-LLM cloud ensemble (\textit{glm-5.1}, \textit{kimi-k2.6}, \textit{deepseek-v4-pro} via Ollama Cloud API) that integrates Layer~1 outputs into a four-tier maintenance criticality label (Critical / High / Medium / Low) per machine via majority-vote aggregation. Pre-registered evaluation against composite ground truth (Lei 2018 + ISO 13374-2) on 25 test machines $\times$ 5 windows = 125 cases yields Quadratic-Weighted Kappa $\kappa = 0.7425$ (95\% CI [0.619, 0.835]) --- substantial agreement per Landis-Koch 1977 --- with perfect Critical-tier precision (1.000) and perfect Critical-or-High coverage of true Critical machines (1.000). Ensemble end-to-end p95 latency of 32.7~s fits within the 60-s Kafka producer cycle. Three pre-registered methodology deviations are documented transparently across V1 ($\kappa = 0.362$), V2 ($\kappa = 0.191$), and V3 ($\kappa = 0.7425$ PASS) iterations.

\textbf{Layer~3} formulates maintenance scheduling as a Quadratic Unconstrained Binary Optimisation (QUBO) over decision variables $x_{m,e,s} \in \{0,1\}$ under industrial-realistic constraints (working hours 09:00--18:00 Mon--Fri, employee groups of 1--20 staff, sequential multi-task, failure-type-specific maintenance durations from CMMS literature). Six solvers are benchmarked across 315 controlled runs and 135 dynamic re-optimisation events: four classical (Greedy, Genetic Algorithm, Simulated Annealing, Tabu Search) all converge to identical optimal makespan with 100\% feasibility and sub-second wall-clock at the largest tested scale (50 machines $\times$ 5 days $\times$ 20 employees = 45{,}000 binary variables). Two quantum-class methods (simulated Quantum Annealing via dwave-neal Ising mode, and quantum-inspired Simulated Bifurcation per Goto et al.\ 2019) consistently fail to produce feasible solutions on this dense scheduling QUBO (paired Wilcoxon $p < 10^{-5}$). Gate-model QAOA on Qiskit Aer statevector simulator is not applicable at meaningful problem sizes due to qubit budget. Real D-Wave hardware was not available for testing. The honest empirical finding at simulator-only evaluation is that classical heuristics dominate quantum-class methods on this problem class; the architectural-readiness positioning for the NISQ-Advanced phase (2025--2028) is anchored on industrial precedents from adjacent scheduling domains (Ford Otosan + D-Wave 6$\times$ scheduling speedup; BASF + D-Wave 7200$\times$ liquid-filling), with the Cerezo et al.\ 2025 quantum-machine-learning skepticism literature explicitly addressed.

End-to-end Layer~2 + Layer~3 pipeline operates within the 60-s real-time Kafka cycle by $\geq 26$~s margin for any classical Layer~3 solver choice. The thesis contributes (1) a validated 3-layer real-time PdM architecture, (2) a multi-LLM cloud ensemble validation methodology with three documented pre-registered deviations and full transparency, and (3) an honest empirical evaluation of classical-versus-quantum scheduling on QASAMAP-realistic problem instances. Future work includes real D-Wave hardware testing, fleet-wide evaluation, and expert-annotation validation of the composite ground truth.

\noindent
\textbf{Keywords:} smart manufacturing, predictive maintenance, anomaly detection, T-GCN spatiotemporal forecasting, agentic AI, large language model ensemble, quantum optimisation, NISQ, scheduling QUBO, real-time Kafka pipeline, pre-registered evaluation
